\section{Availability}
\label{sec:availability}

\begin{center}
\small\sffamily\sloppy
\begin{tabularx}{\linewidth}{@{}L{3.1cm}>{\raggedright\arraybackslash}X@{}}
\ttoprule
\thead{Version} & \ttversion\ (\ttdate)\\
\ttmidrule
\thead{Source} & \url{https://github.com/vlbthambawita/torchtsetlin}\\
\thead{Package} & \url{https://pypi.org/project/torchtsetlin/} --- \code{pip install torchtsetlin}\\
\thead{Documentation} & \url{https://vlbthambawita.github.io/torchtsetlin/}\\
\thead{Licence} & MIT\\
\ttmidrule
\thead{Python} & 3.9, 3.10, 3.11, 3.12\\
\thead{Required} & \code{torch >= 2.0}, \code{numpy >= 1.22} --- no compiler, no CUDA toolkit,
  no C extension\\
\thead{Optional extras} & \code{[viz]} matplotlib \textbar\ \code{[vision]} torchvision
  \textbar\ \code{[sklearn]} scikit-learn \textbar\ \code{[docs]} MkDocs Material
  \textbar\ \code{[dev]} the lot, plus pytest and ruff\\
\ttbotrule
\end{tabularx}
\end{center}

\subsection{Getting started}

\begin{lstlisting}[style=ttshell]
pip install torch          # pick the CPU or CUDA build from pytorch.org
pip install torchtsetlin   # + optional extras: [viz] [vision] [sklearn] [all]
\end{lstlisting}

From source, with the development environment:

\begin{lstlisting}[style=ttshell]
git clone https://github.com/vlbthambawita/torchtsetlin && cd torchtsetlin
pip install -e ".[dev]"
pytest -q                  # the full suite
ruff check src tests       # exactly what CI runs
mkdocs serve               # the documentation site at localhost:8000
\end{lstlisting}

\subsection{Reproducing the measurements}

Every number in Section~\ref{sec:benchmarks} comes from the driver in \code{benchmarks/}. It
runs one suite at a time, merging its JSON after each, so a crash costs only the running
suite, and then regenerates the figures, tables and the documentation page from those records.

\begin{lstlisting}[style=ttshell]
./benchmarks/run_benchmarks.sh                 # every suite, cpu + cuda:0 (~20-25 min)
./benchmarks/run_benchmarks.sh --quick         # ~1 min smoke test of the pipeline
./benchmarks/run_benchmarks.sh -s batch models # a subset
./benchmarks/run_benchmarks.sh --plots-only    # re-render from results already on disk
\end{lstlisting}

The raw records land in \code{benchmarks/results/cpu\_vs\_gpu.json}. Replacing them with your
own hardware's is the whole workflow --- the derived page is regenerated, not hand-edited.

\subsection{Quality gates}

Continuous integration runs \code{ruff} and the full test suite on Python 3.9, 3.11 and 3.12,
then deploys the documentation on pushes to \code{main}. Releases are cut from a \code{v*}
tag through PyPI Trusted Publishing; the release workflow re-runs the suite and smoke-imports
the built wheel in a clean virtual environment before uploading --- a gate added after a
\code{.gitignore} rule silently excluded a subpackage from the 0.1.0 wheel and made
\code{import torchtsetlin} fail for everyone who installed it.

\subsection{Documentation}

The documentation site is MkDocs Material with \code{mkdocstrings} reading Google-style
docstrings directly from the source, so the public API reference cannot drift from the code.
It covers installation and a quick start; concept pages on the algorithm, Booleanization,
batched versus sequential feedback, convolution and hyper-parameters; guides on training,
evaluation, interpretability, visualisation, saving and extending; worked examples; and the
generated CPU-versus-GPU benchmark page.

\subsection{Citation}

\begin{lstlisting}[style=ttshell,basicstyle=\ttfamily\scriptsize\color{ttink}]
@software{torchtsetlin,
  author = {Thambawita, Vajira},
  title  = {torchtsetlin: GPU-enabled, PyTorch-native Tsetlin machines},
  year   = {2026},
  url    = {https://github.com/vlbthambawita/torchtsetlin}
}
\end{lstlisting}

Please also cite the Tsetlin machine papers whose algorithms you rely on --- the reference
list that follows is a reasonable starting point, and the concept pages in the documentation
name the specific paper behind each option.

\subsection{Contributing}

Issues and pull requests are welcome at the repository. The most useful contributions, in
order: a model variant from the literature implemented against the existing hooks
(Section~\ref{sec:architecture}); a fused feedback kernel; validation against the reference
implementations; and applications --- a worked notebook on a real dataset is worth more to the
next user than another synthetic benchmark.
